
\begin{table}[htbp]
\centering
\small
\caption{Established Commitments and Unresolved Questions in the Ontology of Identity}
\label{tab:identity-established-and-unresolved}

\renewcommand{\arraystretch}{1.45}

\begin{tabular}{
    |p{0.43\textwidth}
    |p{0.43\textwidth}|
}
\hline

\textbf{Established Commitments}
&
\textbf{Unresolved Questions}
\\
\hline

Existence and identity are ontologically inseparable.
&
What is the ontological source or origin of a particular identity?
\\
\hline

Neither existence nor identity has temporal or ontological priority over the other.
&
How are particular identities related to the different realms of being?
\\
\hline

Conceptual distinction does not imply ontological separation.
&
Can an existent have multiple identities across different ontological levels?
\\
\hline

An existent is not first constituted as identity-neutral and subsequently endowed with identity.
&
What is the precise relation between identity and the dynamics of being?
\\
\hline

Identity is not reducible to description, classification, recognition, or formal self-equality.
&
Under what conditions can an existent preserve its identity through transformation?
\\
\hline

Determination, distinction, self-sameness, individuation, continuity, and relationality are conceptually distinguishable from identity.
&
What is the ontological status of relations in the constitution or manifestation of identity?
\\
\hline

Epistemic incompleteness does not establish ontological absence or indeterminacy.
&
How should individuation be understood in relation to unity and plurality?
\\
\hline

The identity of an existent need not be fully transparent or exhaustively accessible to a knower.
&
Are there different modes or structures of identity corresponding to different forms of manifestation?
\\
\hline

\end{tabular}

\end{table}
